\documentclass[11pt,a4paper]{article}
\usepackage[utf8]{inputenc}
\usepackage[T1]{fontenc}
\usepackage[french]{babel}
\usepackage[margin=2.4cm]{geometry}
\usepackage{amsmath,amssymb}
\usepackage{booktabs}
\usepackage{xcolor}
\usepackage[colorlinks=true,linkcolor=blue!50!black,citecolor=blue!50!black]{hyperref}
\usepackage{tcolorbox}
\tcbuselibrary{skins,breakable}
\usepackage{titlesec}
\usepackage{enumitem}
\usepackage{parskip}
\setlength{\parskip}{0.4em}

\definecolor{bleuprincip}{HTML}{1F4E78}
\definecolor{bleuclair}{HTML}{2E74B5}
\definecolor{orangealerte}{HTML}{C55A11}
\definecolor{vertok}{HTML}{548235}
\definecolor{violetmethode}{HTML}{7030A0}

\titleformat{\section}{\Large\bfseries\color{bleuprincip}}{\thesection}{0.6em}{}
\titleformat{\subsection}{\large\bfseries\color{bleuclair}}{\thesubsection}{0.6em}{}

\newtcolorbox{intuition}[1][]{colback=bleuclair!7,colframe=bleuclair,boxrule=0.5pt,arc=2pt,
  left=6pt,right=6pt,top=5pt,bottom=5pt,breakable,title={\small\bfseries Intuition},#1}
\newtcolorbox{definitionbox}[1][]{colback=gray!7,colframe=black!60,boxrule=0.4pt,arc=1pt,
  left=6pt,right=6pt,top=5pt,bottom=5pt,breakable,#1}
\newtcolorbox{resultatbox}[1][]{colback=vertok!8,colframe=vertok,boxrule=1pt,arc=2pt,
  left=7pt,right=7pt,top=5pt,bottom=5pt,breakable,#1}
\newtcolorbox{alertebox}[1][]{colback=orangealerte!7,colframe=orangealerte,boxrule=0.6pt,arc=2pt,
  left=7pt,right=7pt,top=5pt,bottom=5pt,breakable,title={\small\bfseries Résultat notable},#1}

\newcommand{\cit}[1]{\textsuperscript{\hyperlink{ref#1}{[#1]}}}

\title{\vspace{-1.3cm}\textcolor{violetmethode}{\Huge\bfseries GDPNow-Maroc}\\[0.3cm]
\textcolor{black}{\Large Méthode 2 --- Rolling AR (modèle de comparaison)}\\[0.15cm]
\textcolor{gray}{\large Réplique du modèle-repère de Higgins (2014), Table 5, modèle 1}}
\author{}
\date{\today}

\begin{document}
\maketitle
\thispagestyle{empty}

\begin{abstract}
\noindent Ce document construit un second modèle, indépendant du premier (\emph{Method\_1\_Higgins}) : un \textbf{Rolling AR} pur, appliqué aux 16 branches, sans aucun indicateur externe ni BVAR --- l'équivalent direct du modèle 1 (« Rolling AR(2) ») de la Table 5 de Higgins\cit{1}, le plus simple de ses 5 modèles de comparaison. Testé sur la même fenêtre hors-échantillon que la Méthode 1 (2021T2--2026T1), il obtient un RMSFE de 14,55 --- nettement moins bon que le modèle complet (5,43), confirmant qualitativement le même écart que Higgins observe dans son propre article.
\end{abstract}

\tableofcontents
\vspace{0.5cm}

% ============================================================================
\section{Pourquoi ce modèle --- sa place dans la comparaison de Higgins}
% ============================================================================

\begin{definitionbox}
Higgins compare son modèle complet (GDPNow) à 5 modèles plus simples (Table 5) :
\begin{center}
\small
\begin{tabular}{lr}
\toprule
\textbf{Modèle (Higgins, USA, 2000--2013)} & \textbf{RMSFE} \\
\midrule
1. Rolling AR(2) & 2,43 \\
2. Rolling Factor Augmented AR(2) & 1,95 \\
3. Mixed Frequency BVAR & 1,73 \\
4. Quarterly BVAR & 2,28 \\
5. Tracking model (mensuel, sans BVAR) & 1,22 \\
6. GDPNow (modèle complet) & \textbf{1,15} \\
\bottomrule
\end{tabular}
\end{center}
\end{definitionbox}

\begin{intuition}
Le modèle 1 de Higgins --- un simple AR(2) roulant, sans aucun indicateur externe --- sert de \textbf{repère minimal} : le modèle complet doit au moins faire mieux que « ne rien savoir d'autre que le passé de la variable elle-même ». C'est exactly ce que cette Méthode 2 réplique pour notre propre travail --- avec une différence assumée et documentée ci-dessous.
\end{intuition}

\subsection{Une différence délibérée avec Higgins}

Higgins fixe l'ordre à $p=2$ pour tous ses AR de comparaison. \textbf{Ici, $p$ est choisi par AIC, indépendamment pour chacune des 16 branches} --- cohérent avec la rigueur déjà appliquée dans notre Méthode 1 (Étape 4), plutôt que de répliquer un choix arbitraire nous-mêmes déjà identifié comme non justifié statistiquement.

% ============================================================================
\section{Méthode}
% ============================================================================

\subsection{L'AR par branche}

Pour chacune des 16 branches, indépendamment :
\begin{equation}
\Delta\log(VA_t) = \alpha + \sum_{k=1}^{p}\gamma_k \Delta\log(VA_{t-k}) + \varepsilon_t, \qquad p^\star = \operatorname*{arg\,min}_{p\in\{1,\dots,8\}} \text{AIC}(p)
\end{equation}
estimé sur le même train (T1-2010--T1-2021) que la Méthode 1, avec la même grille $p\in\{1,\dots,8\}$ que l'Étape 4 du Modèle 1.

\subsection{L'agrégation --- identique à la Méthode 1}

\begin{equation}
\Delta\log(PIB_t) = \sum_{i=1}^{16} SH_{t-1}^i \times \Delta\log(VA_t^{i,\text{AR}})
\end{equation}
Mêmes parts nominales $SH_t^i$, même repli 2010--2013 sur les parts de 2014-T1, exactement la méthode de l'Étape 6 du Modèle 1 --- \textbf{seule la source de $\Delta\log(VA_t^i)$ change} (AR seul, pas de BVAR ni de bridge equation).

\subsection{Le test hors-échantillon --- même fenêtre, pour une comparaison directe}

Les coefficients (dont $p^\star$, jamais retouché) sont appliqués aux 20 mêmes trimestres que la Méthode 1 (2021T2--2026T1), jamais vus par aucune estimation.

% ============================================================================
\section{Résultats}
% ============================================================================

\subsection{Les 16 ordres $p^\star$ retenus}

\begin{center}
\small
\begin{tabular}{lrr}
\toprule
\textbf{Branche} & \textbf{$p^\star$} & \textbf{$R^2$ (train)} \\
\midrule
Autres services & 8 & 0,804 \\
Transports & 8 & 0,482 \\
Administration publique & 4 & 0,450 \\
Éducation-santé & 4 & 0,424 \\
Construction & 6 & 0,381 \\
Pêche & 1 & 0,386 \\
Industrie d'extraction & 5 & 0,308 \\
Industrie de transformation & 4 & 0,322 \\
Services aux entreprises & 4 & 0,246 \\
Électricité, gaz, eau & 4 & 0,207 \\
Immobilier & 4 & 0,194 \\
Hébergement-restauration & 4 & 0,192 \\
Agriculture & 4 & 0,186 \\
Information-communication & 3 & 0,093 \\
Commerce & 1 & 0,048 \\
Finances et assurances & 1 & 0,001 \\
\bottomrule
\end{tabular}
\end{center}

\subsection{Résultat en-échantillon vs hors-échantillon}

\begin{resultatbox}
\begin{center}
\begin{tabular}{lrr}
\toprule
& \textbf{En-échantillon (train)} & \textbf{Hors-échantillon (test)} \\
\midrule
Corrélation & 0,462 & $-0{,}060$ \\
$n$ & 44 & 20 \\
\bottomrule
\end{tabular}
\end{center}
\end{resultatbox}

\begin{alertebox}
\textbf{La corrélation devient négative hors-échantillon} --- le modèle Rolling AR pur, qui semblait capter un signal modeste sur le train (0,462), \textbf{n'a plus aucun pouvoir prédictif réel} une fois confronté à de vraies données jamais vues (RMSFE = 14,55 points de croissance annualisée). Un AR pur, sans aucune information externe à la branche elle-même, s'avère instable et peu fiable pour une vraie prévision.
\end{alertebox}

% ============================================================================
\section{Comparaison complète --- les trois résultats côte à côte}
% ============================================================================

\begin{center}
\small
\begin{tabular}{lrrl}
\toprule
\textbf{Modèle} & \textbf{RMSFE} & \textbf{Corrélation} & \textbf{Contexte} \\
\midrule
Higgins --- Rolling AR(2) & 2,43 & --- & USA, 2000--2013, $p=2$ fixe \\
Higgins --- GDPNow complet & \textbf{1,15} & --- & USA, 2000--2013 \\
\midrule
\textbf{Méthode 2 --- Rolling AR (nous)} & \textbf{14,55} & $-0{,}060$ & Maroc, 2021--2026, $p^\star$ par AIC \\
\textbf{Méthode 1 --- modèle complet (nous)} & \textbf{5,43} & 0,292 & Maroc, 2021--2026 \\
\bottomrule
\end{tabular}
\end{center}

\begin{resultatbox}
\textbf{Le constat qualitatif de Higgins se confirme chez nous} : le modèle complet bat nettement le simple AR --- un rapport de 2,1 pour Higgins (2,43/1,15), un rapport de 2,7 pour nous (14,55/5,43), du même ordre de grandeur relatif. \textbf{En valeur absolue, nos deux modèles sont nettement moins précis que ceux de Higgins} --- cohérent avec les limites déjà identifiées (échantillon de test 3 fois plus petit, indicateurs marocains moins denses que les séries américaines).
\end{resultatbox}

% ============================================================================
\section{Sorties produites}
% ============================================================================

\texttt{resultats/recapitulatif\_AR\_16branches.csv}, \texttt{resultats/PIB\_agrege\_AR.csv}, \texttt{resultats/metriques\_test\_AR.csv}, \texttt{figures/} (17 figures : 16 par branche + 1 agrégat).

\section*{Références}
\addcontentsline{toc}{section}{Références}
\begin{enumerate}[leftmargin=1.6em, label=\textbf{[\arabic*]}]
\item \hypertarget{ref1}{} Higgins, P. (2014). \textit{GDPNow: A Model for GDP ``Nowcasting''}. Federal Reserve Bank of Atlanta, Working Paper 2014-7, Table 5.
\end{enumerate}

\end{document}
